\scanpage{4}
\mspage{3}
(et je suis resté muet sur ce support dans la mouture précédente GF $\underline{\overline{\mathrm{VI}}}$). Quand on interprète $F$ comme figure ensembliste dans l'ens. des lieux $\unsure{L|\mathcal L}$ \deleted{\unsure{ou|de|ce}} (pratiquement cela signifie qu'on a un atelier ensembliste), on l'écrira $\mathrm{mulomb}(F)$ (petite multiombre), son support sera $\mathrm{omb}(F)$.

\uline{Données} \uline{\unsure{pre|pré}\deleted{\unsure{l|li}}mières} :
\[
\left\{\begin{aligned}
&\unsure{\mathscr F|\mathcal F} &&\text{ensemble de figures}\\
&\leq &&\text{\unsure{rellation|relation} d'\uline{inclusion} \unsure{de|des} figures}\\
& &&F\leq G\ \overset{\text{\unsure{déf|def}}}{\Longleftrightarrow}\ F\ \text{sous figure de }G\\
&\ll &&\text{relation de raffinement}\\
& && F\ll G\ \overset{\text{def}}{\Longleftrightarrow}\ F\ \text{raffine }G
\end{aligned}\right.
\]
avec
\[
F\leq G\ \Longrightarrow\ F\ll G\tag*{\makebox[0pt][r]{\makebox[\displaywidth][l]{\uline{At} 0}}}
\]
% slanted
\margin{NB. Ce sera aussi contenu dans \uline{At}~4 plus bas}
-- axiome si évident que je ne l'avais pas inclus dans les axiomes numérotés.

\diagram{\(\begin{array}{@{}clp{.53\linewidth}@{}}
\text{\textcircled{\scriptsize a}} & F\vee G\ \text{ou }F\cup G &
\uline{figure} ``\uline{engendrée}'' ou ``\uline{réunion}'' ou ``\uline{somme}'' : c'est le sup au sens de $\leq$ (quand il existe\unsure{)|\}} -- auquel cas c'est aussi un sup pour $\ll$\\
\end{array}\)}
\margin{\deleted{\unsure{Question}}\inserted{NB.} : Si $(F,G)$ ont un sup pour $\ll$ (c'est tjrs le cas dans le cas ensembliste $\mathrm{Fig}(L)$), ce n'est pas néc. un sup pour $\leq$}
\margin{\uline{NB} $\widetilde{F\cap G}=\widetilde F\cap\widetilde G$}
\diagram{\(\begin{array}{@{}clp{.53\linewidth}@{}}
\text{\textcircled{\scriptsize b}} & \begin{aligned}&F\wedge G\ \text{ou } F\cap G\\&\quad(=\mathrm{Inf}^{\leq}(F,\unsure{c|G}))\end{aligned}
& figure \uline{intersection}\inserted{ou Inf pour $\leq$} -- existe toujours \unsure{(|et (}$\overline{\mathrm m}$ pour familles infinie (non vide)) ; quand en $=$ temps $F\vee C$ existe, alors c'est aussi le Inf pour $\ll$.\\
\end{array}\)}
\margin{\uline{NB} Existe dans le cas ensembliste $\mathrm{Fig}(L)$ et dans le cas \unsure{mechisé|méchisé}}
\diagram{\(\begin{array}{@{}clp{.53\linewidth}@{}}
\text{\textcircled{\scriptsize\unsure{c}}} & F\mathbin{\ldotp}G\quad (=\mathrm{Inf}^{\ll}(F,c)) &
figure ``\uline{intersection} \uline{fine} (quand elle existe) : Inf pour $\ll$.\\
&\deleted{\unsure{G|CG}}&\\
\text{\textcircled{\scriptsize\unsure{d}}} & \varnothing_{\mathscr F} &
figure vide : plus petit élément de $\mathscr F$ (pour $\leq$ \uline{et} pour $\ll$)\\
\end{array}\)}
